\section{Discussion}
\label{discussion}

%%%%%%%%%%%%%%%%%%%%%%%%%%%%%
% SECTION 5.1 BELOW
%%%%%%%%%%%%%%%%%%%%%%%%%%%%%

\subsection{Contribution to an EWS-oriented FEM modelling workflow}
\label{discussion_ews_fem_workflow}

This project developed a progressively more complete COMSOL finite-element route for simulating dynamic breast deformation in the context of the Early Warning Scan (EWS). The main contribution was not a clinically validated detection model, but a reproducible model-development framework in which anatomical structure, material parameters, support assumptions, loading route and tumor-overlay placement can be varied in a controlled way. This distinction is important because EWS is based on observable breast-surface motion, whereas many of the model quantities used during development, such as internal stress or volume-averaged displacement, are primarily mechanical diagnostics.

Compared with a simplified breast representation, the final model route introduced a curved posterior support geometry, a more realistic glandular distribution, optional volumetric skin representation, material-parameter sensitivity cases, Cooper-like support approximations, dynamic loading variants and analytic tumor-overlay placements. These additions increased the anatomical and mechanical richness of the model and made it possible to evaluate which components most strongly affect the simulated surface response. The framework therefore provides a basis for future EWS-oriented surface-motion studies, but the present results should be interpreted as model-development outcomes rather than as evidence of diagnostic performance.


%%%%%%%%%%%%%%%%%%%%%%%%%%%%%
% SECTION 5.2 BELOW
%%%%%%%%%%%%%%%%%%%%%%%%%%%%%

\subsection{Anatomical realism and glandular representation}
\label{discussion_anatomical_realism_glandular}

A central aim of this project was to improve the anatomical realism of the FEM breast model. The realistic glandular route added a spatially distributed fibroglandular structure instead of representing the internal breast tissue with a single simplified region. This is relevant because breast deformation depends not only on total breast size, but also on the relative amount and spatial distribution of adipose and fibroglandular tissue \cite{Samani2007,Ramiao2016,Goodbrake2022,Teixeira2023}.

The current glandular implementation also introduced a practical limitation. A detailed glandular structure is mechanically more representative, but it increases geometric complexity, meshing difficulty and solver cost. Similar anatomical complexity was a limiting factor in previous internal FEBio-based model-development routes \cite{EngelsEWSReport,StormEWSReport}. The current COMSOL workflow was able to build and solve the realistic glandular route, but running many dynamic variants on a local laptop remained computationally demanding. This suggests that the modelling strategy is technically feasible, but that further studies require access to stronger computational resources.

Anatomical variability should therefore become a larger part of future model development. Breast size and shape, glandular fraction, glandular asymmetry, posterior attachment geometry and tissue distribution vary strongly between individuals \cite{Ramiao2016,Goodbrake2022,Teixeira2023}. The current model can be used as a first step for building such population-level model sets, but this would require many additional variants and more validation. Future work should therefore vary the relevant anatomical, mechanical and loading parameters in a systematic way to quantify their influence on the simulated surface response. For future AI-based analysis, this type of systematic variation could support the construction of synthetic datasets with controlled differences in anatomy, tissue properties, loading conditions and tumor characteristics.

%%%%%%%%%%%%%%%%%%%%%%%%%%%%%
% SECTION 5.3 BELOW
%%%%%%%%%%%%%%%%%%%%%%%%%%%%%

\subsection{Material-parameter uncertainty and response realism}
\label{discussion_material_uncertainty_response_realism}

Material parameters remain one of the largest sources of uncertainty in breast FEM modelling. Reported stiffness values for adipose tissue, fibroglandular tissue, skin and tumor tissue vary substantially across studies because they depend strongly on experimental protocol, tissue condition and constitutive assumptions \cite{Samani2007,Ramiao2016,Goodbrake2022,Teixeira2023}. The material settings used in this project should therefore be interpreted as controlled reference and sensitivity values rather than as patient-specific tissue properties.

This uncertainty affects how the displacement results should be read. A simulated displacement amplitude is only physically meaningful when the material parameters, boundary conditions and loading route are plausible for the intended experimental situation. The current results show that the model can produce measurable dynamic surface motion, but the exact displacement scale still requires comparison with experimental or EWS-like motion data. Such validation would make it possible to determine whether the selected stiffness range reproduces realistic breast motion.

The von Mises stress results should be interpreted with caution. In this project, von Mises stress was useful as a scalar diagnostic for comparing stress concentrations between variants, especially near support and attachment regions. However, stress is not directly measured by the EWS concept and should not be treated as a direct detection output. Stress concentrations near the posterior support are more likely to reflect the imposed boundary condition and support geometry than a clinically meaningful tissue response. For EWS interpretation, surface displacement and landmark motion are therefore more relevant than internal stress metrics.

%%%%%%%%%%%%%%%%%%%%%%%%%%%%%
% SECTION 5.4 BELOW
%%%%%%%%%%%%%%%%%%%%%%%%%%%%%

\subsection{Dynamic loading and surface-motion interpretation}
\label{discussion_dynamic_loading_surface_motion}

The dynamic loading definition is an important part of the model, because the simulated breast response depends strongly on how motion is introduced. The fixed-support acceleration input and the prescribed support-displacement route represent different idealisations of how breast motion could be generated during an EWS-like measurement. The fixed-support acceleration route applies a controlled global excitation while keeping the support geometry fixed. The prescribed support-displacement route is more directly related to a physical situation in which the chestwall or support moves up and down, for example during a jump-like motion.

In the current project, most quantitative dynamic comparisons focused on the fixed-support acceleration route. This provided a controlled way to compare model variants, but it does not necessarily mean that this route is the most realistic representation of the final EWS setup. Further work should therefore compare the fixed-support acceleration route with prescribed support-motion approaches in more detail. The resulting motion depends on the amplitude, timing, damping, gravity preload and support coupling of the model. These factors should therefore be selected or calibrated against the intended EWS motion setting, because they directly influence whether the simulated breast response is physically representative.

The current dynamic results should therefore be interpreted as proof that different loading routes can be implemented and compared within the COMSOL workflow. They do not yet define an experimentally validated EWS loading protocol. A future validation study should record the motion applied in an EWS-like setup and use that measured motion as the input or calibration target for the FEM model.

%%%%%%%%%%%%%%%%%%%%%%%%%%%%%
% SECTION 5.5 BELOW
%%%%%%%%%%%%%%%%%%%%%%%%%%%%%

\subsection{Skin and Cooper-like support modelling choices}
\label{discussion_skin_cooper_support}

The skin representation had a clear mechanical role in the model because it added a relatively stiff outer layer around the softer internal breast tissues. This is consistent with multi-component breast FEM studies, where the breast is represented using separate anatomical and mechanical components rather than as one homogeneous soft-tissue volume \cite{Chen2025,Teixeira2023,Mazier2024}. In the current COMSOL route, the $1.5~\mathrm{mm}$ volumetric skin layer was therefore an important step toward a more anatomically representative model. This thickness is also within reported breast-skin thickness ranges, although skin thickness and mechanical properties vary between subjects and measurement locations \cite{Huang2008,Sutradhar2012}.

The volumetric skin route should still be interpreted as a controlled modelling choice rather than as a patient-specific skin reconstruction. Its mechanical effect depends not only on the selected skin stiffness and thickness, but also on how the layer is coupled to the internal breast tissue. Future work should therefore include a more systematic skin-parameter study to determine how robust the simulated surface response is across realistic skin properties.

The Cooper-like support implementation should also be interpreted as an approximation. Cooper's ligaments form a distributed fibrous support structure in the breast, while the current COMSOL route represented their influence through simplified support terms rather than explicit anatomical ligament reconstruction \cite{Chen2025}. This made it possible to test the mechanical influence of additional support, but future versions should use a more anatomically support representation for the Cooper's ligaments.

%%%%%%%%%%%%%%%%%%%%%%%%%%%%%
% SECTION 5.6 BELOW
%%%%%%%%%%%%%%%%%%%%%%%%%%%%%

\subsection{Tumor-overlay route and future tumor sensitivity}
\label{discussion_tumor_overlay_future_sensitivity}

The tumor-overlay route successfully added a build-level mechanism for placing analytic spherical stiffness inclusions inside the existing breast volume. This is a useful first step because tumor size and location can be changed without rebuilding the full breast geometry for every variant. The current implementation is therefore suitable for preparing systematic tumor-sensitivity studies.

The current tumor representation is still simplified. Real breast tumors can differ in morphology, stiffness, depth and interaction with surrounding tissue, and reported tumor stiffness values vary strongly between measurement methods and tumor types \cite{Samani2007,McKnight2002,Patel2021,Patel2022}. A spherical homogeneous overlay can test the effect of a controlled stiffness contrast, but it cannot represent the full range of malignant tumor morphology. This limitation is especially important for EWS, because a tumor would only be detectable through its effect on the externally observable surface motion. The relevant question is therefore not only whether a stiff inclusion changes the internal mechanical field, but whether it changes the surface response enough to be measured robustly.

Future tumor studies should therefore test a broader range of tumor geometries, stiffness contrasts and anatomical locations. These studies should be run as complete dynamic simulations with tumor-mask post-processing, surface-motion metrics and matched no-tumor controls. Only after such a solved and validated case set is available can the tumor-overlay route be used to discuss potential EWS detectability.

%%%%%%%%%%%%%%%%%%%%%%%%%%%%%
% SECTION 5.7 BELOW
%%%%%%%%%%%%%%%%%%%%%%%%%%%%%

\subsection{Computational limitations and reproducibility}
\label{discussion_computational_limitations_reproducibility}

A major limitation of the current project was that all COMSOL model builds and time-dependent simulations were run on a personal laptop. This was workable for developing the pipeline and generating geometry build-only cases, but it strongly limited the number of complete dynamic solves that could be included. For the more realistic model variants, individual dynamic simulations could take approximately $12~\mathrm{h}$ per case, not including the additional time required for post-processing. This made large parameter sweeps and complete dynamic testing of the more complex model routes impractical within the project timeframe.

Access to faster external computational resources would improve the study by allowing larger sweeps, more realistic model variants and full post-processing to be run more efficiently. This would be especially important for systematic testing of material parameters, dynamic loading routes, Cooper-like support and tumor-overlay cases. The current computational limitation should therefore be interpreted as a practical limitation of the project setup and timeframe, rather than as a limitation of the modelling framework itself. Future work should account for this by running the simulations on stronger computational hardware.

The repository structure helped keep the model-development workflow traceable. The scripted case definitions, compact outputs, report notes and figure sources made it possible to track how the model variants were generated and evaluated \cite{KuijpersEWSComsolGitHub}. Future work should keep this level of traceability, especially when larger simulation sets are generated on external computational resources.


%%%%%%%%%%%%%%%%%%%%%%%%%%%%%
% SECTION 5.8 BELOW
%%%%%%%%%%%%%%%%%%%%%%%%%%%%%

\subsection{Overall perspective and future work}
\label{discussion_overall_perspective_future_work}

Overall, the project extended the COMSOL FEM breast model in several directions that are relevant for future EWS-oriented deformation studies:

\begin{itemize}
    \item a more anatomically detailed breast geometry with chestwall curvature, asymmetry options and a realistic glandular tissue distribution;
    \item controlled material-parameter and volumetric-skin variants;
    \item simplified Cooper-like support and dynamic loading routes;
    \item analytic tumor-overlay placement for future stiffness-sensitivity studies;
    \item a traceable workflow for case definition, model building and automatic post-processing.
\end{itemize}

The most important next step is to run systematic, fully solved dynamic simulations on stronger computational hardware and compare the simulated surface response with experimental or EWS-like measurement data. This should be combined with controlled variation of material parameters, glandular anatomy, skin properties, loading route and tumor properties. If those future simulations show that tumor-related stiffness contrasts produce measurable and repeatable differences in breast-surface motion, the model could become a useful tool for exploring EWS sensitivity and design choices.

